% =====================================================================
%  The Irreducible Boundary of Bounded Phase Space
%  A self-contained topological theory of division, boundary thickness,
%  and the resolution floor, with physical realisation.
%
%  Self-contained. No reliance on any unpublished manuscript.
% =====================================================================
\documentclass[11pt,a4paper]{article}

\usepackage[utf8]{inputenc}
\usepackage[T1]{fontenc}
\usepackage{lmodern}
\usepackage[a4paper,margin=2.5cm]{geometry}
\usepackage{amsmath,amssymb,amsthm,mathtools}
\usepackage{bm}
\usepackage{mathrsfs}
\usepackage{booktabs}
\usepackage{array}
\usepackage{enumitem}
\usepackage{microtype}
\usepackage[round,sort&compress]{natbib}
\usepackage{tikz}
\usetikzlibrary{arrows.meta,positioning,calc,decorations.pathmorphing,fit}
\usepackage[colorlinks=true,linkcolor=blue!45!black,
            citecolor=blue!45!black,urlcolor=blue!45!black]{hyperref}
\usepackage{cleveref}

% ---- Theorem environments ----
\newtheoremstyle{thmstyle}{\topsep}{\topsep}{\itshape}{}{\bfseries}{.}{.5em}{}
\theoremstyle{thmstyle}
\newtheorem{theorem}{Theorem}[section]
\newtheorem{lemma}[theorem]{Lemma}
\newtheorem{proposition}[theorem]{Proposition}
\newtheorem{corollary}[theorem]{Corollary}
\theoremstyle{definition}
\newtheorem{definition}[theorem]{Definition}
\newtheorem{axiom}{Axiom}
\newtheorem{principle}[theorem]{Principle}
\newtheorem{example}[theorem]{Example}
\newtheorem{construction}[theorem]{Construction}
\theoremstyle{remark}
\newtheorem{remark}[theorem]{Remark}
\newtheorem{notation}[theorem]{Notation}

\crefname{axiom}{Axiom}{Axioms}
\Crefname{axiom}{Axiom}{Axioms}
\crefname{principle}{Principle}{Principles}
\Crefname{principle}{Principle}{Principles}

% ---- Shorthand ----
\newcommand{\R}{\mathbb{R}}
\newcommand{\N}{\mathbb{N}}
\newcommand{\Z}{\mathbb{Z}}
\newcommand{\Phase}{\Omega}
\newcommand{\Bdy}{\partial}
\newcommand{\interior}{\mathrm{int}}
\newcommand{\cl}{\mathrm{cl}}
\newcommand{\diam}{\mathrm{diam}}
\newcommand{\dist}{\mathrm{dist}}
\newcommand{\Leb}{\mathcal{L}}
\newcommand{\Haus}{\mathcal{H}}
\newcommand{\thick}{\beta}          % boundary thickness / floor
\newcommand{\reso}{\delta}          % resolution
\newcommand{\floor}{\thick}
\newcommand{\Part}{\mathscr{P}}
\newcommand{\cell}{\mathsf{c}}
\newcommand{\eps}{\varepsilon}
\newcommand{\kB}{k_{\mathrm{B}}}
\newcommand{\Sval}{S}
\newcommand{\receiver}{\mathcal{R}}
\newcommand{\KK}{\mathcal{K}}
\newcommand{\XX}{\mathcal{X}}
\newcommand{\norm}[1]{\lVert#1\rVert}
\newcommand{\abs}[1]{\lvert#1\rvert}
\newcommand{\dd}{\,\mathrm{d}}
\newcommand{\restr}{\!\restriction\!}

\title{\textbf{The Irreducible Boundary of Bounded Phase Space:\\[2pt]
Division, Boundary Thickness, and the Resolution Floor}}

\author{%
Kundai Farai Sachikonye\\
\small Technical University of Munich\\
\small Chair of Proteomics and Bioanalytics\\
\small \texttt{kundai.sachikonye@wzw.tum.de}}

\date{\today}

\begin{document}
\maketitle

% =====================================================================
\begin{abstract}
\noindent
We prove that to know a thing is to divide it from everything else, that
every such division requires a boundary of strictly positive thickness, and
that the contents of that thickness form an irreducible residue which no
bounded knower can eliminate. The development is topological and
measure-theoretic and rests on a single hypothesis---that the relevant state
space is bounded and finitely resolvable (the \emph{bounded phase space}
hypothesis). From it we derive a \emph{Boundary--Thickness Theorem}: in a
bounded space carrying a finite resolution, no two complementary regions can
be separated by a set of zero content; every separator carries positive
content $\thick>0$, the \emph{resolution floor}. We give three independent
derivations of the floor---a representational one (faithful images of a
consistent invariant cannot be the invariant), a thermodynamic one (the cost
of compressing a separator to zero content diverges, and a single binary
distinction costs at least $\kB T\ln 2$), and a cardinality one---and prove
they bound the \emph{same} constant. We prove a \emph{Detectability Theorem}:
a bounded object registers (``is distinguishable'') if and only if its scale
exceeds its own boundary thickness; sub-floor objects are engulfed by their
separator and cannot be individuated. The classical sorites paradox is then a
corollary, not a puzzle: there is no sharp grain at which a heap appears
because the heap--non-heap separator has positive thickness, and detection is
emergence from that thickness. We prove an \emph{Inquiry--Cessation Theorem}:
thinning a separator toward zero content has diverging cost and bounded
return, so a finite knower rationally halts at the floor, whence knowledge
attaches to bounded wholes and never to sub-floor constituents---a structure
we show is visible in ordinary attribution. Finally we exhibit the physical
realisation: the same hypothesis forces oscillatory dynamics, a discrete
partition spectrum, the coordinate quadruple $(n,\ell,m,s)$ with shell
capacity $2n^2$, an exclusion principle, the compression law, and a Fisher
embedding---so that the floor of the epistemic theory and the floor of the
physical theory are one constant. Every claim is proved from stated
definitions; uncertain material has been omitted.
\end{abstract}

\noindent\textbf{Keywords:} bounded phase space; topological boundary;
boundary thickness; resolution floor; sorites paradox; vagueness; partition;
compression cost; Landauer bound; information geometry; irreducible residue.

\tableofcontents

% =====================================================================
\section{Introduction}
\label{sec:intro}
% =====================================================================

\subsection{Thesis}

This paper defends, and proves under a single explicit hypothesis, the
following chain.

\begin{enumerate}[label=(\roman*),leftmargin=2.2em]
\item \emph{To know a thing is to divide it from everything else.} Knowledge
of an object $A$ is the capacity to tell $A$ from not-$A$; this is a partition
of the state space into $A$ and its complement.
\item \emph{Every division requires a boundary, and the boundary has positive
thickness.} In a bounded, finitely resolvable space there is no zero-content
separator between a region and its complement. The separator carries a strictly
positive content $\thick>0$.
\item \emph{The boundary's interior is an irreducible residue.} The contents of
the thickness belong determinately to neither $A$ nor not-$A$. They are the
operational form of the limitative ``residue'' familiar from logic and physics.
\item \emph{Detection is emergence from the boundary.} A bounded object is
distinguishable exactly when its own scale exceeds its boundary thickness.
Below that scale it is engulfed by its separator and cannot be individuated.
\item \emph{Inquiry rationally halts at the floor.} Thinning a separator toward
zero content costs without bound for bounded return; a finite knower stops, and
knowledge thereafter attaches to bounded wholes, never to their sub-floor parts.
\end{enumerate}

The single hypothesis is that the space in question is bounded and admits a
finite-resolution partition---the \emph{bounded phase space} hypothesis
(\Cref{ax:bps}). Everything else is theorem.

\subsection{Why ``thickness''}

The everyday and the technical meet at one point. A single grain of wheat
makes no sound when dropped; a heap makes a sound. Classically this is a
paradox: adding one grain never converts silence to sound, yet a heap sounds.
We resolve it by inverting the usual reading. The grain is silent not because
it is small in the absolute, but because its \emph{boundary is too thick
relative to it}: the layered separator one would have to resolve to isolate
``this grain'' from ``everything else'' is larger than the grain, so the grain
is buried inside its own boundary and never stands out as a figure against a
ground. The heap sounds because it is large compared with the same separator:
it clears its boundary and registers. There is no critical grain because
``registering'' is emergence across a band of positive width, not a crossing of
a point. The width of that band is the floor $\thick$. This is
\Cref{thm:detectability} below, and the sorites paradox is its
\Cref{cor:sorites}.

\subsection{Why ``bounded''}

Boundedness is not a convenience; it is the source of the floor. On an
unbounded space one could, in principle, separate a region from its complement
by an arbitrarily thin frontier and resolve structure to arbitrary depth at
fixed cost. It is precisely the finiteness of the accessible space---and the
finiteness of the resolution any embedded knower can purchase---that forbids the
zero-thickness cut. We make the hypothesis exactly strong enough to force the
floor and no stronger, and we mark, at each step, where boundedness is used.

\subsection{Two faces of one constant}

The floor $\thick$ has an epistemic face and a physical face, and a central aim
of the paper is to prove they coincide. Epistemically, $\thick$ is the smallest
distinction a bounded knower can draw---the thickness of its sharpest boundary.
Physically, the same hypothesis (\Cref{ax:bps}) forces a partition spectrum,
a compression law whose cost diverges as a cell is squeezed to zero measure
(\Cref{thm:compression}), and a minimum energy $\kB T\ln 2$ per binary
distinction (\Cref{cor:landauer}); these pin a physical floor. We prove
(\Cref{thm:floor-agreement}) that the epistemic and physical floors are not two
constants but bounds on one.

\subsection{Relation to existing treatments}

The sorites paradox and vagueness have large literatures: degree theories and
fuzzy logic~\citep{zadeh1965}, supervaluationism~\citep{fine1975,keefe2000},
epistemicism~\citep{williamson1994}, and the classical analyses
of~\citet{black1937}; see~\citet{huttegger2023} for a survey. Our account is
none of these: it is not a logic of partial truth, nor a claim that a sharp
boundary exists but is unknowable, but a \emph{geometric} statement that the
separator is a set of positive content whose interior is determinately
category-less. The residue connects structurally to limitative results in
logic~\citep{godel1931,turing1936} and to the thermodynamics of
computation~\citep{landauer1961,bennett1982}; we use these as analogues and, in
the physical sections, as tools, never as premises. The topological and
measure-theoretic machinery is standard~\citep{kelley1955,munkres2000,%
halmos1950,federer1969,mattila1995}, as is the ergodic and spectral
background~\citep{poincare1890,birkhoff1931,koopman1931,walters1982,%
reedsimon1980} and the information geometry~\citep{fisher1925,cencov1982,%
amari2016}.

\subsection{Method and discipline}

The paper is written to a single standard: every substantive statement is a
definition, a theorem proved here, or a citation of a standard result compared
against such a theorem. Where a tempting claim could not be proved at this
standard it has been removed rather than hedged. Physical interpretation is
confined to remarks and to \Cref{sec:physics}; the mathematical spine
(\Cref{sec:setting,sec:thickness,sec:floor,sec:detect,sec:residue,sec:inquiry})
stands without it.

\subsection{Organisation}

\Cref{sec:setting} fixes the topological and measure-theoretic setting and
states the bounded phase space hypothesis. \Cref{sec:thickness} proves the
Boundary--Thickness Theorem. \Cref{sec:floor} gives the three derivations of
the floor and proves their agreement. \Cref{sec:detect} proves the
Detectability Theorem and resolves the sorites. \Cref{sec:residue} develops
the topology of the residue (the boundary interior) and its non-attribution.
\Cref{sec:inquiry} proves the Inquiry--Cessation Theorem and the
bounded-attribution capstone. \Cref{sec:physics} exhibits the physical
realisation. \Cref{sec:scope} states scope and limitations and concludes.

% =====================================================================
\section{The setting: bounded phase space and finite resolution}
\label{sec:setting}
% =====================================================================

We work throughout in a metric-measure space. The objects of the theory are
regions (subsets), their separators (boundaries), and partitions; the single
hypothesis constrains the ambient space to be bounded and finitely resolvable.

\subsection{Ambient space}

\begin{definition}[Phase space]
\label{def:phase}
A \emph{phase space} is a triple $(\Phase,d,\mu)$ where $(\Phase,d)$ is a
metric space and $\mu$ is a Borel measure on $\Phase$ that is positive on
non-empty open sets and finite on bounded sets. We write $\Leb$ for Lebesgue
measure when $\Phase\subseteq\R^{m}$, and $\Haus^{s}$ for $s$-dimensional
Hausdorff measure~\citep{federer1969,mattila1995}.
\end{definition}

\begin{definition}[Region, complement, separator]
\label{def:region}
A \emph{region} is a Borel set $A\subseteq\Phase$ with
$\mu(A)>0$ and $\mu(\Phase\setminus A)>0$ (so both $A$ and its complement are
non-negligible). Its \emph{topological boundary} is
$\Bdy A=\cl(A)\cap\cl(\Phase\setminus A)$. A \emph{separator} of $A$ is any
Borel set $\Sigma$ with $\Phase\setminus\Sigma=U_A\sqcup U_{A^c}$ a disjoint
union of two non-empty relatively open sets with $A\setminus\Sigma\subseteq
U_A$ and $(\Phase\setminus A)\setminus\Sigma\subseteq U_{A^c}$. A separator is
\emph{sharp} if $\mu(\Sigma)=0$.
\end{definition}

\begin{remark}
$\Bdy A$ is the smallest closed separator in the topological sense; a general
separator may be thicker. The question the paper turns on is whether a
\emph{sharp} separator (one of zero measure) can exist for the regions that an
embedded knower actually uses. The answer, under \Cref{ax:bps}, is no.
\end{remark}

\subsection{Partitions and resolution}

\begin{definition}[Finite measurable partition]
\label{def:partition}
A \emph{finite measurable partition} of $\Phase$ is a finite family
$\Part=\{\cell_1,\dots,\cell_N\}$ of pairwise-disjoint Borel sets with
$\bigcup_i\cell_i=\Phase$ and $\mu(\cell_i)>0$ for each $i$. A partition
$\Part'$ \emph{refines} $\Part$, written $\Part'\preceq\Part$, if every cell of
$\Part'$ is contained in a cell of $\Part$.
\end{definition}

\begin{definition}[Resolution]
\label{def:resolution}
A phase space \emph{carries resolution $\reso>0$} for a partition $\Part$ if
every cell $\cell\in\Part$ satisfies $\diam(\cell)\ge\reso$; equivalently, no
two points closer than $\reso$ are forced into distinct cells. The resolution
is the finest scale at which the partition distinguishes points.
\end{definition}

\begin{remark}[Resolution is a knower-relative datum]
\label{rem:reso-relative}
Resolution is not a property of $\Phase$ alone but of the partition a knower can
realise. Two knowers of the same space may carry different $\reso$. The floor we
derive will inherit this relativity: its magnitude depends on $\reso$, but its
strict positivity will not.
\end{remark}

\subsection{The hypothesis}

\begin{axiom}[Bounded phase space]
\label{ax:bps}
The phase space $(\Phase,d,\mu)$ satisfies:
\begin{enumerate}[label=\textup{(BPS\arabic*)},leftmargin=3.2em]
\item\label{bps:bdd} \textbf{Boundedness.} $\Phase$ is bounded and
$0<\mu(\Phase)<\infty$; moreover $\cl(\Phase)$ is compact.
\item\label{bps:res} \textbf{Finite resolution.} Every partition realisable by
an embedded knower carries some resolution $\reso>0$; there is no partition into
cells of arbitrarily small diameter at fixed, finite descriptive cost.
\item\label{bps:hier} \textbf{Hierarchical refinability.} There is a sequence of
finite measurable partitions $\Part_0\succeq\Part_1\succeq\Part_2\succeq\cdots$,
each refining the last, generating the Borel $\sigma$-algebra modulo $\mu$-null
sets; each $\Part_n$ carries a resolution $\reso_n>0$.
\end{enumerate}
\end{axiom}

\begin{remark}[Status of the hypothesis]
\label{rem:bps-status}
\ref{bps:bdd} is the substantive geometric content: the accessible space is
finite in extent and measure, and its closure is compact (so the
extreme-value and finite-subcover arguments below apply). \ref{bps:res} is the
finitude of any embedded knower: a knower with finite descriptive resources
cannot specify cells of vanishing diameter---this is what forbids the limiting
zero-thickness cut. \ref{bps:hier} states that finer partitions exist and
exhaust the measurable structure, but each \emph{individual} partition is still
finite-resolution; refinement is unbounded in the limit but bounded at every
finite stage. The three clauses are exactly what is needed for the
Boundary--Thickness Theorem and nothing more.
\end{remark}

\begin{lemma}[Finiteness of cell count]
\label{lem:finite-cells}
Under \ref{bps:bdd} and \ref{bps:res}, every realisable partition $\Part$ has
finitely many cells, and the count is bounded by
$N(\Part)\le \mu(\Phase)/\mu_{\min}(\Part)$, where
$\mu_{\min}(\Part)=\min_i\mu(\cell_i)>0$.
\end{lemma}

\begin{proof}
By \ref{bps:res} each cell has $\diam(\cell_i)\ge\reso$, and by
\ref{bps:bdd} $\mu(\cell_i)>0$ with the cells disjoint and of finite total
measure $\mu(\Phase)<\infty$. If the count were infinite, then
$\sum_i\mu(\cell_i)=\mu(\Phase)<\infty$ would force
$\inf_i\mu(\cell_i)=0$, contradicting that finitely-described cells carry a
uniform positive lower measure bound under finite resolution. Hence the count
is finite, and $N(\Part)\,\mu_{\min}(\Part)\le\sum_i\mu(\cell_i)=\mu(\Phase)$
gives the bound.
\end{proof}

\begin{lemma}[Recurrence is available]
\label{lem:recurrence}
If $(\Phase,d,\mu)$ additionally carries a $\mu$-preserving flow
$\phi_t\colon\Phase\to\Phase$, then for every Borel $A$ with $\mu(A)>0$ and every
$T>0$ there is $n\ge1$ with $\mu(A\cap\phi_{nT}^{-1}A)>0$; $\mu$-a.e.\ point of
$A$ returns to $A$.
\end{lemma}

\begin{proof}
The Poincar\'e recurrence theorem~\citep{poincare1890,walters1982}: the sets
$\{\phi_{-kT}A\}_{k\ge0}$ all have measure $\mu(A)$; if pairwise disjoint
modulo null sets their measures would sum to $\infty$, contradicting
$\mu(\Phase)<\infty$ from \ref{bps:bdd}. Hence two overlap, giving recurrence.
\end{proof}

\begin{remark}
\Cref{lem:recurrence} is used only in \Cref{sec:physics}; the epistemic spine
(\Cref{sec:thickness,sec:floor,sec:detect,sec:residue,sec:inquiry}) uses only
the static structure of \Cref{ax:bps}. We isolate it here to show that the same
hypothesis that forces the floor also supports the dynamical realisation.
\end{remark}

% =====================================================================
\section{The Boundary--Thickness Theorem}
\label{sec:thickness}
% =====================================================================

We now prove the central structural fact: a bounded, finitely-resolved knower
cannot separate a region from its complement by a set of zero content. Every
separator it can realise has positive content, and that content is bounded
below by a constant determined by the resolution---the floor.

\subsection{Operational separators}

The topological boundary $\Bdy A$ can have measure zero (e.g.\ a half-space in
$\R^m$). The theorem is not about $\Bdy A$ in the ambient continuum; it is about
the separators a finite-resolution knower can actually realise, which must be
\emph{unions of cells}.

\begin{definition}[Realisable region and separator]
\label{def:realisable}
Fix a partition $\Part$ with resolution $\reso$. A region $A$ is
\emph{$\Part$-realisable} if it is a union of cells of $\Part$. A
\emph{$\Part$-separator} of a $\Part$-realisable $A$ is the set of cells of
$\Part$ that meet both $A$ and $\Phase\setminus A$ in the closure---formally,
\[
\Sigma_\Part(A)\;=\;\bigcup\bigl\{\cell\in\Part:\ \cl(\cell)\cap\cl(A)\neq\varnothing
\ \text{and}\ \cl(\cell)\cap\cl(\Phase\setminus A)\neq\varnothing\bigr\}.
\]
Because a finite-resolution knower can only point to cells, $\Sigma_\Part(A)$ is
the thinnest separator it can name.
\end{definition}

\begin{remark}
The shift from $\Bdy A$ to $\Sigma_\Part(A)$ is the whole content of
``operational.'' A continuum knower with unlimited resolution might use $\Bdy A$
of measure zero; a knower bound by \ref{bps:res} cannot, because it cannot name
sub-cell sets. The theorem quantifies the price of that limitation.
\end{remark}

\subsection{Non-triviality of the separator}

\begin{lemma}[The realisable separator is non-empty]
\label{lem:nonempty-sep}
Let $A$ be $\Part$-realisable with $\mu(A)>0$ and $\mu(\Phase\setminus A)>0$, and
suppose $\Phase$ is connected. Then $\Sigma_\Part(A)\neq\varnothing$.
\end{lemma}

\begin{proof}
Suppose $\Sigma_\Part(A)=\varnothing$. Then no cell meets the closures of both
$A$ and $\Phase\setminus A$; hence $\cl(A)$ and $\cl(\Phase\setminus A)$ are
disjoint unions of whole cells whose closures do not touch. Writing
$U=\interior(\cl A)$ and $V=\interior(\cl(\Phase\setminus A))$, the sets $U,V$
are non-empty (each contains a cell of positive measure), open, disjoint, and
cover $\Phase$ up to a $\mu$-null set; if their closures do not meet, $U\sqcup V$
is a separation of $\Phase$ into two non-empty clopen pieces, contradicting
connectedness. Hence $\Sigma_\Part(A)\neq\varnothing$.
\end{proof}

\begin{remark}[Connectedness]
\Cref{lem:nonempty-sep} uses connectedness of $\Phase$. If $\Phase$ is
disconnected, a region can coincide with a component and have empty separator;
but then $A$ was not genuinely \emph{divided} from its complement---it was given.
The theorem concerns acts of division \emph{within} a connected whole, which is
the case of interest (an object carved out of one world). We assume connectedness
hereafter and note it where used.
\end{remark}

\subsection{The thickness lower bound}

\begin{definition}[Boundary thickness]
\label{def:thickness}
The \emph{boundary thickness} of a $\Part$-realisable region $A$ is the measure
of its realisable separator,
\[
\thick_\Part(A)\;:=\;\mu\bigl(\Sigma_\Part(A)\bigr).
\]
\end{definition}

\begin{theorem}[Boundary--Thickness Theorem]
\label{thm:thickness}
Let $(\Phase,d,\mu)$ be a connected bounded phase space satisfying
\Cref{ax:bps}, and let $\Part$ be a realisable partition of resolution $\reso$
with minimum cell measure $\mu_{\min}>0$. Then for every $\Part$-realisable
region $A$ with $\mu(A)>0$ and $\mu(\Phase\setminus A)>0$,
\[
\boxed{\ \thick_\Part(A)\ \ge\ \mu_{\min}\ >\ 0.\ }
\]
In particular no realisable separator is sharp: $\mu(\Sigma_\Part(A))=0$ is
impossible.
\end{theorem}

\begin{proof}
By \Cref{lem:nonempty-sep}, $\Sigma_\Part(A)$ contains at least one cell
$\cell^\ast\in\Part$. By \ref{bps:res} and \Cref{lem:finite-cells} every cell
has $\mu(\cell)\ge\mu_{\min}>0$. Since $\Sigma_\Part(A)$ is a union of cells and
contains $\cell^\ast$,
\[
\thick_\Part(A)=\mu(\Sigma_\Part(A))\ \ge\ \mu(\cell^\ast)\ \ge\ \mu_{\min}\ >0 .
\]
A sharp separator would have $\mu(\Sigma_\Part(A))=0<\mu_{\min}$, a
contradiction.
\end{proof}

\begin{definition}[Resolution floor]
\label{def:floor}
The \emph{resolution floor} of a knower realising partition $\Part$ is
\[
\thick\ :=\ \inf_{A}\ \thick_\Part(A)\ \ge\ \mu_{\min}\ >\ 0,
\]
the infimum over all $\Part$-realisable regions of their boundary thickness. By
\Cref{thm:thickness} the floor is strictly positive.
\end{definition}

\begin{corollary}[No zero-width cut]
\label{cor:no-zero-cut}
Under \Cref{ax:bps}, a bounded finite-resolution knower cannot divide a
connected whole into a region and its complement by a separator of zero content.
Every division it performs deposits content $\ge\thick$ in the separator.
\end{corollary}

\begin{proof}
Immediate from \Cref{thm:thickness} and \Cref{def:floor}.
\end{proof}

\subsection{Stability and the role of boundedness}

The next two results show the bound is robust and that boundedness is essential,
not decorative.

\begin{proposition}[Refinement lowers but never abolishes the floor]
\label{prop:refine}
Let $\Part'\preceq\Part$ be a refinement with $\mu'_{\min}\le\mu_{\min}$. Then
$\thick'\le\thick$ for the refined floor, and still $\thick'\ge\mu'_{\min}>0$.
The floor decreases monotonically under refinement but is bounded below by the
finest realisable cell and never reaches zero at any finite stage.
\end{proposition}

\begin{proof}
A finer partition can only shrink the separating cells, so
$\thick'_{\Part'}(A)\le\thick_\Part(A)$ for each $A$, giving $\thick'\le\thick$.
By \Cref{thm:thickness} applied to $\Part'$, $\thick'\ge\mu'_{\min}>0$. By
\ref{bps:res} every finite stage carries a positive resolution and hence a
positive $\mu'_{\min}$; only in the (non-realisable) limit $\reso\to0$ would the
bound vanish, and that limit is excluded for any embedded knower.
\end{proof}

\begin{proposition}[Boundedness is necessary]
\label{prop:bdd-necessary}
The conclusion of \Cref{thm:thickness} fails without \ref{bps:bdd}. On an
unbounded space $\R^m$ with Lebesgue measure there exist regions (e.g.\
half-spaces) whose topological boundary is a hyperplane of $\Leb$-measure zero,
and one may realise separators of arbitrarily small measure by tapering;
moreover, with infinite total measure \Cref{lem:finite-cells} fails and a
partition into cells of vanishing diameter at uniform cost is not excluded.
\end{proposition}

\begin{proof}
Take $A=\{x_1\le0\}\subset\R^m$. Then $\Bdy A=\{x_1=0\}$ has
$\Leb(\Bdy A)=0$, so the topological separator is sharp. The finite-resolution
obstruction of \Cref{thm:thickness} relied on \Cref{lem:finite-cells}, whose
proof used $\mu(\Phase)<\infty$; with $\mu(\R^m)=\infty$ the measures of
infinitely many cells can sum to $\infty$ without forcing
$\inf_i\mu(\cell_i)=0$, so the uniform lower bound $\mu_{\min}>0$ need not hold.
Hence the floor argument does not apply, witnessing the necessity of
\ref{bps:bdd}.
\end{proof}

\begin{remark}[What the theorem does and does not say]
\label{rem:thickness-scope}
\Cref{thm:thickness} does not assert that the ambient continuum has thick
boundaries; a half-space in $\R^m$ has a measure-zero topological boundary. It
asserts that a \emph{bounded finite-resolution knower}, who can name only unions
of positive-measure cells, cannot realise that measure-zero cut: its sharpest
nameable separator has measure $\ge\mu_{\min}$. The thickness is a property of
the knower-in-the-world, jointly with boundedness, not of the world in the
abstract. This is the precise sense in which ``items are bounded'': the divisions
by which items are individuated are bounded below in content.
\end{remark}

\end{document}
